\chapter{Allergies}

\section*{Definition and Overview}
An allergy is an abnormal immune response to a normally harmless substance, known as an allergen. In a susceptible person, the immune system mistakenly identifies the allergen as a threat and mounts an inflammatory reaction that produces symptoms such as sneezing, itching, rashes, wheezing, or digestive upset. Common allergens include pollens, house dust mites, animal dander, moulds, foods, insect venom, and some medicines. Reactions range from mild seasonal discomfort to severe, potentially life-threatening anaphylaxis. Allergies are among the most common chronic conditions in the world, can begin at any age, and frequently run in families. They are distinct from intolerances, such as lactose intolerance, which do not involve the immune system and are generally less dangerous.

\section*{Epidemiology}
Allergic disorders are estimated to affect around one in four people in many developed countries, and their prevalence has risen substantially over recent decades, particularly in children. The pattern of allergy changes with age: food allergy and eczema are common in infancy and early childhood, while allergic rhinitis and asthma frequently emerge in the school years and adolescence. A person with one allergic condition is at increased risk of developing others, a phenomenon known as the allergic (atopic) march. This tendency to become sensitised is strongly influenced by genetics, and a family history of allergy is the most consistent risk factor. Environmental factors, including urban living, air pollution, dietary patterns, and the timing of allergen exposure, also contribute to the rising prevalence.

\section*{Aetiology and Pathophysiology}
Allergies develop through a process of sensitisation in genetically predisposed individuals. On first exposure, allergen-presenting cells drive the production of allergen-specific immunoglobulin E (IgE), which attaches to mast cells and basophils. On subsequent exposure, the allergen cross-links the bound IgE, causing the release of histamine and other inflammatory mediators. These chemicals produce the immediate symptoms of allergy, and a later cellular phase maintains inflammation over hours to days.

Common allergens include:

\begin{itemize}
    \item \textbf{Inhaled:} pollens, house dust mites, mould spores, and animal dander.
    \item \textbf{Ingested:} foods such as egg, cow's milk, peanut, tree nuts, wheat, soy, and seafood.
    \item \textbf{Injected:} insect venom from bee, wasp, and ant stings, and some medicines.
    \item \textbf{Contact:} latex, nickel, and some cosmetics or topical preparations.
    \item \textbf{Drugs:} antibiotics such as penicillins, non-steroidal anti-inflammatory drugs, and anaesthetic agents.
\end{itemize}

The site and nature of the reaction depend on the route of exposure, the amount of allergen, and the individual's pattern of sensitisation.

\section*{Clinical Features}
Symptoms depend on how the allergen enters the body and where the reaction occurs:

\begin{itemize}
    \item \textbf{Nose and eyes:} sneezing, runny or blocked nose, and itchy, watery eyes in allergic rhinitis and conjunctivitis.
    \item \textbf{Skin:} urticaria (hives), itching, red patches, and flare-ups of eczema.
    \item \textbf{Chest:} wheeze, cough, and breathlessness in allergic asthma.
    \item \textbf{Gut:} vomiting, abdominal pain, and diarrhoea after ingestion of food allergens.
    \item \textbf{Mouth:} itching or swelling of the lips, tongue, and throat after eating the offending food.
    \item \textbf{General:} malaise and fatigue during prolonged allergic reactions.
    \item \textbf{Severe (anaphylaxis):} difficulty breathing, throat swelling, dizziness, and collapse, which can be fatal.
\end{itemize}

Symptoms typically appear within minutes to hours of exposure and often follow a recognisable pattern on repeated exposure to the same allergen.

\section*{Differential Diagnosis}
Conditions that can be mistaken for allergy include:

\begin{itemize}
    \item \textbf{Food intolerances:} lactose intolerance, food sensitivities, and reactions to additives that do not involve the immune system.
    \item \textbf{Viral infections:} colds and other viral illnesses that mimic hay fever or skin rashes.
    \item \textbf{Non-allergic rhinitis:} triggered by smoke, temperature change, or stress without an allergic basis.
    \item \textbf{Gastrointestinal disorders:} irritable bowel syndrome and other gut conditions that cause symptoms after food.
    \item \textbf{Irritant skin reactions:} from chemicals, soaps, or heat rather than immune-mediated allergy.
    \item \textbf{Drug side effects:} adverse effects of medicines that resemble allergic reactions.
    \item \textbf{Anxiety and hyperventilation:} which can produce tingling, breathlessness, and dizziness similar to a mild reaction.
\end{itemize}

\section*{When to Seek Medical Care}
See a doctor if allergic symptoms are frequent, severe, or interfering with sleep, school, work, or growth, and before any food group is removed from a child's diet, because unsupervised elimination can cause nutritional problems. People who have had a severe reaction should seek specialist assessment for an adrenaline auto-injector and a written action plan.

\textbf{Call 000 (or local emergency services) for:}
\begin{itemize}
    \item Difficulty breathing, wheezing, or a hoarse voice after allergen exposure.
    \item Swelling of the tongue, face, or throat, or difficulty swallowing.
    \item Dizziness, collapse, or loss of consciousness.
    \item Any symptoms of anaphylaxis, even if they begin mildly.
\end{itemize}

\section*{Diagnosis and Investigations}
Diagnosis combines a careful history with targeted tests:

\begin{itemize}
    \item \textbf{Detailed history:} what was eaten or encountered, the timing of symptoms, the speed of onset, and the response to treatment or avoidance.
    \item \textbf{Clinical examination:} identification of rashes, urticaria, or signs of allergic rhinitis or asthma.
    \item \textbf{Skin prick testing:} small amounts of allergen are placed on the skin to look for a local wheal, indicating sensitisation.
    \item \textbf{Serum specific IgE:} blood tests measuring allergen-specific antibodies.
    \item \textbf{Food diaries and elimination diets:} under professional supervision to link symptoms to specific foods.
    \item \textbf{Oral food or sting challenges:} supervised exposure in a clinical setting for selected cases where the diagnosis is uncertain.
    \item \textbf{Referral:} to an allergy specialist or clinical immunologist for complex, severe, or unclear presentations.
\end{itemize}

\section*{Management}
There is no cure for allergy, but most allergic conditions can be well controlled:

\begin{itemize}
    \item \textbf{Allergen avoidance:} identifying and avoiding the specific allergen where practical.
    \item \textbf{Antihistamines:} oral or topical preparations for sneezing, itching, and hives.
    \item \textbf{Intranasal corticosteroids and inhalers:} for allergic rhinitis and allergic asthma.
    \item \textbf{Topical steroids and emollients:} for eczema.
    \item \textbf{Adrenaline auto-injectors:} prescribed for people at risk of anaphylaxis, with training in their use.
    \item \textbf{Allergen immunotherapy:} gradual desensitisation that reduces sensitivity to certain allergens, such as pollens, dust mites, and insect venom.
    \item \textbf{Written action plans:} clear instructions for recognition and treatment of reactions, shared with family, school, and workplace.
\end{itemize}

\section*{Complications}
Complications of allergy include:

\begin{itemize}
    \item Severe reactions (anaphylaxis) from accidental exposure.
    \item Poorly controlled allergic asthma leading to emergency department attendance.
    \item Chronic sinusitis and nasal polyps from untreated rhinitis.
    \item Growth or nutritional problems if foods are avoided unnecessarily.
    \item Sleep disturbance, poor concentration, and reduced quality of life.
    \item Anxiety about eating, being stung, or being exposed to allergens.
\end{itemize}

\section*{Prognosis and Outcomes}
Many allergies improve with age. Children commonly outgrow allergies to cow's milk, egg, and wheat, whereas allergies to peanut, tree nuts, seafood, and insect venom are more likely to persist. The majority of people achieve good control of symptoms with avoidance and appropriate medication. Anaphylaxis is uncommon but is a serious risk that requires an emergency action plan and, where indicated, an adrenaline auto-injector. For selected allergens, immunotherapy offers lasting improvement that can continue for years after a completed course.

\section*{Prevention and Self-Care}
Simple measures reduce the impact of allergies:

\begin{itemize}
    \item Know your triggers and carry a written allergy action plan.
    \item Read food and product labels carefully.
    \item Use allergen-proof bedding covers if allergic to dust mites, and wash bedding in hot water.
    \item Check pollen counts and adjust outdoor activity during allergy season.
    \item Keep an adrenaline auto-injector accessible, check its expiry, and rehearse its use.
    \item Inform family, school, and employers about your allergy.
    \item Speak to a doctor before changing a child's diet to manage allergies.
\end{itemize}

\section*{Sources and Further Reading}
The following reputable sources provide reliable, up-to-date information on allergies:

\begin{itemize}
    \item NHS. \textit{Allergies information for patients.} https://www.nhs.uk/
    \item Mayo Clinic. \textit{Allergies: symptoms, causes, and treatment.} https://www.mayoclinic.org/
    \item MSD Manual. \textit{Overview of allergic and atopic disorders.} https://www.msdmanuals.com/
    \item MedlinePlus. \textit{Allergy health information.} https://medlineplus.gov/
    \item NICE. \textit{Clinical guidance on allergy and allergic conditions.} https://www.nice.org.uk/
    \item World Health Organization. \textit{Information on allergic diseases and prevention.} https://www.who.int/
\end{itemize}
